\section{Introduction}\label{sec:intro}

Reference-conditioned identity preservation -- feed a reference portrait,
get the same character in a new scene -- is normally bought one of two
ways: train a dedicated adapter on a large paired dataset
\citep{ye2023ipadapter,li2023photomaker,wang2024instantid}, or switch to
an instruction-following edit model whose base was already pretrained
for image-to-image tasks \citep{blackforest2025fluxkontext}. Neither
route is available when the base model must remain a frozen, pure
text-to-image (T2I) diffusion transformer: no per-subject or general
fine-tuning, no departure from T2I semantics into edit-model territory,
no library modification. We show such a model -- a frozen anime T2I DiT,
Z-Image \citep{zimage2025} -- can be activated as a reference-conditioned
identity generator by two cheap, mechanistically independent
interventions: concatenating the reference image's attention key/value
tensors into a restricted band of self-attention layers at inference
time (zero training), and a rank-8 LoRA on 0.06\% of backbone parameters
within the same layer band (minutes of training, not per-subject).

Getting either mechanism to actually work exposes the same blind spot
twice. Both mechanisms have a natural, cheap-to-measure headline metric
-- does the generated character look like the reference? -- and both
times, optimizing or ablating against that metric alone would have led
to the wrong conclusion. The training-free mechanism's naive
position-preserving key injection causes literal copy-paste ``ghosting''
because the reference and target token grids share the same RoPE
coordinate origin; a frequency-gated fix removes the ghosting and *does*
raise raw identity-match rate on a small development set (23\%$\to$40\%)
-- but held-out validation with a defect-aware judge shows the same
recipe also raises a severe composition-collapse defect rate from 20\%
to 82\%, so the rate of an actually usable image collapses from 53\% to
15\%. The training-light mechanism's layer-position ablation was
pre-registered to test whether a specific 16-layer band is uniquely
responsible for identity-conditioning capacity; identity-match turns out
to be statistically flat across every tested band (95\% CI overlapping,
confirmed under subject-level clustering, and replicated on a second
held-out subject set) -- the hypothesis fails. An apparent 8.6$\times$
defect-rate gradient across the same bands then failed our own audit:
it does not replicate cross-set, and removing a generation-protocol
confound from its reference arm erases it on the original set.

\paragraph{Contributions.}
\textbf{(1) A mechanistic fix and its honest limit.} A mechanistic
explanation, from the transformer's own RoPE coordinate construction, of
why naive position-preserving key/value injection produces copy-paste
ghosting, and a frequency-gated re-application of RoPE that removes it --
together with held-out evidence that the fix,
as currently tuned, trades ghosting for a worse composition-collapse
failure mode, so we recommend the simpler position-free baseline for
deployment rather than the mechanism this project set out to promote.
\textbf{(2) A training-light alternative, characterized by a
pre-registered null result.} A 0.06\%-parameter band-LoRA competitive
with a purpose-trained 12.9B adapter baseline on identity-match and
usable-output rate, whose designed keystone ablation (does layer
position determine identity capacity?) returns a null that replicates
across two held-out subject sets -- and whose one apparent positive
side-finding we refuted ourselves, by cross-set replication and a
confound-free re-run, and report as an audited artifact.
\textbf{(3) A cross-mechanism evaluation finding.} Across two
independent mechanisms, the decisive practical differences lived on
axes standard evaluation ignores: a composition-defect axis invisible
to identity-match scoring (training-free line), and a
subject-set/generation-protocol sensitivity that manufactured a
CI-decisive but spurious ablation effect (training-light line). Future
reference-conditioning work should report a joint identity-and-defect
metric, replicate across subject sets, and hold generation protocol
fixed across ablation arms.
\textbf{(4) A held-out validation protocol} that catches exactly the
failures it was built to catch, on probe sets disjoint from every
development set used to tune either mechanism -- including refuting
this paper's own strongest intermediate finding.
